\documentclass[11pt,a4paper]{article}

\usepackage[utf8]{inputenc}
\usepackage[T1]{fontenc}
\usepackage{lmodern}
\usepackage{float}
\usepackage[margin=2cm]{geometry}
\usepackage{amsmath,amssymb}
\usepackage{booktabs}
\usepackage{graphicx}
\usepackage{hyperref}
\usepackage{xcolor}
\usepackage{titlesec}
\usepackage{tabularx}

\hypersetup{
    colorlinks=true,
    linkcolor=blue!70!black,
    citecolor=blue!70!black,
    urlcolor=blue!70!black
}

\titleformat{\section}
{\large\bfseries}
{\thesection.}
{0.5em}
{}

\titleformat{\subsection}
{\normalsize\bfseries}
{\thesubsection.}
{0.5em}
{}

\titlespacing*{\section}
{0pt}
{1.2ex plus 0.5ex minus .2ex}
{0.6ex}

\titlespacing*{\subsection}
{0pt}
{1.0ex plus 0.4ex minus .2ex}
{0.4ex}

\setlength{\parskip}{0.45em}
\setlength{\parindent}{0pt}

\begin{document}

\begin{center}
{\LARGE \textbf{Projeto 1: Regressão com MLP e Estudo de Ablação}} \\[0.35em]
{\large Disciplina: Introdução às Redes Neurais Artificiais} \\[0.3em]
\textbf{Autor:} Gabriely Di Folco Rocha (RA: 188892) \\[0.2em]
\textbf{Repositório:} \url{https://github.com/ghostfoxgoddes3/ANN}
\end{center}

\vspace{-0.5em}
\hrule
\vspace{0.8em}


\section{Configuração Experimental e Reprodutibilidade}

\subsection{Inspeção e Caracterização dos Dados}

O conjunto de dados utilizado no projeto, denominado
\texttt{dataset\_p1}, possui 300 observações e duas variáveis numéricas,
$x$ e $y$. A variável $x$ foi utilizada como atributo de entrada da
rede, enquanto $y$ corresponde à variável-alvo da tarefa de regressão.

Não foram identificados valores ausentes no conjunto de dados. As duas
variáveis apresentam tipo numérico de ponto flutuante
(\texttt{float64}). A variável $x$ possui 300 valores distintos,
distribuídos no intervalo $[0,10]$, com uma observação para cada valor.
Além disso, os valores de $x$ são igualmente espaçados ao longo do
intervalo. A variável $y$ apresenta valores entre aproximadamente
$-1,49$ e $2,66$.

A Tabela~\ref{tab:dataset_summary} apresenta as principais estatísticas
descritivas das duas variáveis.

\begin{table}[H]
\centering
\small
\caption{Estatísticas descritivas do conjunto de dados.}
\label{tab:dataset_summary}
\begin{tabular}{lrr}
\toprule
\textbf{Estatística} & \textbf{$x$} & \textbf{$y$} \\
\midrule
Contagem & 300 & 300 \\
Média & 5.000000 & 0.465705 \\
Desvio padrão & 2.901229 & 0.726711 \\
Mínimo & 0.000000 & -1.488851 \\
1º quartil & 2.500000 & -0.021091 \\
Mediana & 5.000000 & 0.452895 \\
3º quartil & 7.500000 & 0.988865 \\
Máximo & 10.000000 & 2.655266 \\
\bottomrule
\end{tabular}
\end{table}

A inspeção gráfica indica que a relação entre $x$ e $y$ não é linear,
apresentando comportamento oscilatório e dispersão das observações ao
longo do domínio de $x$. Essa característica justifica a utilização de
um modelo com capacidade de representar relações não lineares. A
Figura~\ref{fig:dataset} apresenta a distribuição das observações.

\begin{figure}[H]
\centering
\includegraphics[width=0.82\textwidth]{figures/dataset.png}
\caption{Distribuição das observações do conjunto de dados.}
\label{fig:dataset}
\end{figure}


\subsection{Pré-processamento e Partição dos Dados}

O \texttt{dataset\_p1} foi dividido segundo duas estratégias distintas,
ambas mantendo a proporção de 10\% para treinamento, 10\% para
validação e 80\% para teste. Como o conjunto possui 300 observações,
cada estratégia resultou em 30 amostras para treinamento, 30 para
validação e 240 para teste.

A primeira estratégia consiste em uma partição aleatória, utilizando
uma semente fixa igual a 42 para garantir a reprodutibilidade. A segunda
consiste em uma partição sistemática baseada nos índices das
observações.

Na partição sistemática, os índices foram classificados de acordo com
o resto da divisão por 30. Os índices cujo resto pertence ao conjunto
$\{0,10,20\}$ foram destinados ao treinamento, enquanto aqueles cujo
resto pertence a $\{5,15,25\}$ foram destinados à validação. As demais
observações foram destinadas ao teste.

Como os valores de $x$ são igualmente espaçados e a classificação é
baseada nos índices, essa construção distribui as observações de
treinamento e validação ao longo de todo o domínio de $x$. Assim, a
partição sistemática não concentra as amostras de treinamento em uma
única região do conjunto de dados.

As duas estratégias foram mantidas ao longo de todo o experimento. Isso
permite verificar se os efeitos das modificações introduzidas na MLP
dependem da composição dos subconjuntos ou se permanecem semelhantes
sob diferentes estratégias de particionamento.

Para o pré-processamento, foi utilizado o método
\texttt{StandardScaler}. Os parâmetros de transformação foram ajustados
exclusivamente sobre os dados de treinamento, evitando a utilização de
informações dos conjuntos de validação e teste nessa etapa. Para uma
variável genérica $v$, a transformação é dada por

\[
z =
\frac{v-\mu_{\mathrm{train}}}
{\sigma_{\mathrm{train}}},
\]

onde $\mu_{\mathrm{train}}$ e $\sigma_{\mathrm{train}}$ representam,
respectivamente, a média e o desvio-padrão calculados sobre o conjunto
de treinamento.

Os parâmetros obtidos no treinamento foram então utilizados para
transformar os conjuntos de validação e teste. O procedimento foi
aplicado separadamente à variável de entrada $x$ e à variável resposta
$y$, utilizando-se um escalonador específico para cada variável. Dessa
forma, as informações das partições de validação e teste não participam
do ajuste dos parâmetros de padronização.

Durante a avaliação final, as previsões foram convertidas novamente
para a escala original de $y$ por meio da transformação inversa. Assim,
as métricas utilizadas durante o treinamento e a seleção puderam ser
calculadas no espaço padronizado, enquanto os resultados finais também
puderam ser apresentados na unidade original da variável resposta.


\subsection{Definição e Otimização Empírica do Modelo Baseline (Vanilla)}

O modelo \textit{baseline} foi definido como uma MLP básica treinada
estritamente com \textbf{SGD padrão}, sem aceleradores adaptativos como
Adam e sem termo de \textit{momentum}. Também não foi utilizada
regularização explícita por L1, L2 ou Dropout.

A busca empírica pela configuração do modelo envolveu a variação do
número de camadas, do número de neurônios por camada, da taxa de
aprendizado inicial ($\eta$), da função de ativação e do tamanho do lote
(\textit{batch size}). O critério de seleção foi exclusivamente o menor
MSE observado no conjunto de validação.

A configuração de referência proposta para a MLP foi:

\[
\text{Entrada}
\rightarrow
\text{Dense}(128)
\rightarrow
\text{ReLU}
\rightarrow
\text{Dense}(64)
\rightarrow
\text{ReLU}
\rightarrow
\text{Saída}(1).
\]

Os hiperparâmetros utilizados na busca são apresentados na
Tabela~\ref{tab:baseline_hiper}. A busca foi realizada
separadamente para as partições aleatória e sistemática.

\begin{table}[H]
\centering
\small
\caption{Espaço de busca e configurações selecionadas para o modelo
Baseline (Vanilla).}
\label{tab:baseline_hiper}
\begin{tabularx}{\textwidth}{lXl}
\toprule
\textbf{Hiperparâmetro} &
\textbf{Valores explorados} &
\textbf{Configuração escolhida} \\
\midrule
Topologia das camadas &
$(32)$, $(64)$, $(128)$, $(32,32)$, $(64,64)$ e $(128,64)$ &
$(128,64)$ para ambas \\

Função de ativação &
ReLU e Tanh &
ReLU para ambas \\

Otimizador &
SGD padrão, sem momentum &
SGD, momentum = 0 \\

Taxa de aprendizado &
$0,001$ e $0,01$ &
$0,01$ para ambas \\

Tamanho do lote &
8, 16 e 32 &
16 (aleatória); 32 (sistemática) \\

Função de custo &
Erro Quadrático Médio (MSE) &
MSELoss \\

Critério de parada &
\textit{Early Stopping}, monitorando o erro de validação &
Patience = 20 épocas \\

Semente aleatória &
Semente fixa para reprodutibilidade &
Seed = 42 \\
\bottomrule
\end{tabularx}
\end{table}

O espaço de busca correspondeu a 72 combinações de hiperparâmetros para
cada estratégia de partição. A seleção foi realizada exclusivamente a
partir do menor MSE obtido no conjunto de validação, mantendo o conjunto
de teste intocado.

Para a partição aleatória, a configuração selecionada foi composta por
duas camadas ocultas, com 128 e 64 neurônios, respectivamente, ativação
ReLU, taxa de aprendizado $0,01$ e tamanho de lote 16. O menor MSE de
validação foi $0,818838$, obtido na época 564. O treinamento foi
encerrado na época 584 após o critério de \textit{Early Stopping}.

Para a partição sistemática, a configuração selecionada também foi
composta por duas camadas ocultas, com 128 e 64 neurônios, ativação ReLU
e taxa de aprendizado $0,01$, mas com tamanho de lote 32. O menor MSE de
validação foi $0,770725$, obtido na época 998. Nesse caso, o treinamento
atingiu o limite máximo de 1000 épocas.

Os resultados resumidos da busca são apresentados na
Tabela~\ref{tab:baseline_resultados}.

\begin{table}[H]
\centering
\small
\caption{Melhores resultados obtidos na busca do modelo Baseline
(Vanilla).}
\label{tab:baseline_resultados}
\begin{tabular}{lrrrrr}
\toprule
\textbf{Partição} &
\textbf{Arquitetura} &
\textbf{Ativação} &
\textbf{Batch} &
\textbf{MSE val.} &
\textbf{Melhor época} \\
\midrule
Aleatória &
$(128,64)$ &
ReLU &
16 &
0,818838 &
564 \\

Sistemática &
$(128,64)$ &
ReLU &
32 &
0,770725 &
998 \\
\bottomrule
\end{tabular}
\end{table}

A configuração selecionada em cada partição foi então utilizada no
treinamento do modelo baseline, cuja dinâmica de otimização foi
analisada por meio das curvas de MSE e MAE.


\section{Dinâmica de Treinamento e Convergência}

O treinamento dos modelos baseline foi acompanhado por meio do MSE e
do MAE, calculados ao final de cada época nos conjuntos de treinamento
e validação. As duas estratégias de particionamento foram mantidas na
análise.

Além do menor valor do erro de validação, foi registrada a época em que
esse valor ocorreu e o número total de épocas executadas. O estado do
modelo correspondente ao menor MSE de validação foi preservado para as
etapas posteriores.


\subsection{Divisão aleatória}

Para a divisão aleatória, o modelo selecionado na etapa de busca
hiperparamétrica possui arquitetura $(128,64)$, ativação ReLU, taxa de
aprendizado $0,01$, \textit{batch size} 16 e SGD sem momentum.

\begin{figure}[H]
\centering
\includegraphics[width=\textwidth]{figures/baseline_random_curves.png}
\caption{Evolução do MSE e do MAE no treinamento e na validação para
o modelo baseline com divisão aleatória.}
\label{fig:baseline_random_curves}
\end{figure}

O menor MSE de validação foi $0,818838$, obtido na época 564. Com
critério de parada antecipada de 20 épocas sem melhoria, o treinamento
foi encerrado na época 584.

Ao final do treinamento, o MSE no conjunto de treinamento foi
$0,651541$, enquanto o MSE de validação foi $0,873208$. Os respectivos
MAEs foram $0,661778$ e $0,794267$. Esses valores estão na escala
padronizada da variável resposta.

A diferença entre o menor MSE de validação e o valor registrado ao final
do treinamento também mostra a utilidade do \textit{Early Stopping}:
embora o treinamento tenha continuado após a época 564 para confirmar
a ausência de novas melhorias, o modelo selecionado corresponde ao
estado de menor erro de validação.


\subsection{Divisão sistemática}

Para a divisão sistemática, foi utilizada a mesma arquitetura
$(128,64)$, com ativação ReLU, taxa de aprendizado $0,01$ e SGD sem
momentum. O \textit{batch size} selecionado foi 32.

\begin{figure}[H]
\centering
\includegraphics[width=\textwidth]{figures/baseline_systematic_curves.png}
\caption{Evolução do MSE e do MAE no treinamento e na validação para
o modelo baseline com divisão sistemática.}
\label{fig:baseline_systematic_curves}
\end{figure}

Nesse caso, o menor MSE de validação foi $0,770725$, obtido na época
998. Como o critério de parada antecipada não foi atingido antes do
limite máximo, o treinamento foi executado por 1000 épocas.

Ao final do treinamento, o MSE no conjunto de treinamento foi
$0,670215$, enquanto o MAE foi $0,648229$. O MAE de validação
correspondente ao melhor ponto foi $0,734492$.


\subsection{Comparação da dinâmica de treinamento}

Os dois particionamentos apresentaram dinâmicas distintas. Na divisão
aleatória, o menor erro de validação ocorreu na época 564, seguido pela
interrupção na época 584. Na divisão sistemática, o menor erro de
validação ocorreu apenas na época 998, próxima ao limite máximo de
treinamento.

Essa diferença mostra que o número de épocas necessário para atingir o
melhor ponto depende da composição dos dados utilizados no treinamento
e na validação. Também evidencia que velocidade de convergência não
deve ser usada como substituta de uma métrica de generalização: um
modelo pode atingir um ponto estável rapidamente e ainda apresentar
erro de validação elevado.

As curvas de MSE e MAE constituem, portanto, a referência para avaliar
as alterações introduzidas nas etapas de ablação.


\section{Estudo de Ablação e Desempenho no Conjunto de Teste}

Após a definição do modelo baseline, foram avaliadas quatro variações
individuais e uma configuração combinada, seguindo o roteiro
experimental: inclusão de momentum, regularização L2, regularização L1,
Dropout e, por fim, combinação de componentes selecionados a partir dos
experimentos individuais.

Em cada ablação, foram preservadas a arquitetura e as condições do
baseline correspondente ao particionamento, alterando-se o componente
específico sob investigação. Dessa forma, as diferenças observadas
podem ser atribuídas ao mecanismo introduzido dentro das condições do
experimento.

O conjunto de teste, correspondente a 80\% das observações, permaneceu
intocado durante todas as etapas de treinamento e seleção.

O principal critério de comparação foi o menor MSE observado no conjunto
de validação. A MAE foi utilizada como métrica complementar. Os
resultados desta etapa estão no espaço padronizado da variável resposta.


\subsection{Baseline + Momentum}

A primeira ablação consistiu na inclusão de momentum no otimizador SGD,
utilizando $\beta=0,9$, sem alteração da arquitetura, taxa de
aprendizado ou tamanho do lote.

\begin{table}[H]
\centering
\small
\caption{Resultados da ablação com momentum.}
\label{tab:momentum_ablation}
\begin{tabular}{lrrrr}
\toprule
\textbf{Split} &
\textbf{MSE val.} &
\textbf{MAE val.} &
\textbf{Melhor época} &
\textbf{Épocas} \\
\midrule
Random &
0,701855 &
0,689057 &
152 &
172 \\

Systematic &
0,920469 &
0,769003 &
7 &
27 \\
\bottomrule
\end{tabular}
\end{table}

No particionamento aleatório, a introdução do momentum reduziu o MSE
de validação de $0,818838$ para $0,701855$, uma redução de
aproximadamente 14,29\%. A MAE também diminuiu, de $0,794267$ para
$0,689057$. O menor erro de validação foi atingido na época 152.

No particionamento sistemático, o comportamento foi diferente. O MSE
de validação aumentou de $0,770725$ para $0,920469$, enquanto a MAE
passou de $0,734492$ para $0,769003$. O melhor ponto ocorreu já na
época 7.

Portanto, dentro das configurações avaliadas, o efeito do momentum não
foi uniforme entre os dois particionamentos. A alteração do mecanismo
de atualização modificou tanto o desempenho de validação quanto a
dinâmica de convergência.

\begin{figure}[H]
\centering
\includegraphics[width=0.90\textwidth]{figures/momentum_random_curves.png}
\caption{Dinâmica de treinamento do modelo com momentum no
particionamento aleatório.}
\label{fig:momentum_random_curves}
\end{figure}

\begin{figure}[H]
\centering
\includegraphics[width=0.90\textwidth]{figures/momentum_systematic_curves.png}
\caption{Dinâmica de treinamento do modelo com momentum no
particionamento sistemático.}
\label{fig:momentum_systematic_curves}
\end{figure}


\subsection{Baseline + Regularização L2}

A segunda ablação introduziu regularização L2 por meio do parâmetro
\texttt{weight\_decay}, utilizando $\lambda_{L2}=10^{-4}$ e mantendo o
restante da configuração igual ao baseline.

\begin{table}[H]
\centering
\small
\caption{Resultados da ablação com regularização L2.}
\label{tab:l2_ablation}
\begin{tabular}{lrrrr}
\toprule
\textbf{Split} &
\textbf{MSE val.} &
\textbf{MAE val.} &
\textbf{Melhor época} &
\textbf{Épocas} \\
\midrule
Random &
0,804006 &
0,747255 &
623 &
643 \\

Systematic &
0,770923 &
0,734615 &
999 &
1000 \\
\bottomrule
\end{tabular}
\end{table}

No particionamento aleatório, o L2 reduziu o MSE de validação de
$0,818838$ para $0,804006$, aproximadamente 1,81\%. A MAE também foi
reduzida de $0,794267$ para $0,747255$.

No particionamento sistemático, a alteração foi muito pequena: o MSE
passou de $0,770725$ para $0,770923$ e a MAE passou de $0,734492$ para
$0,734615$.

Assim, para o valor de $\lambda_{L2}$ utilizado, a regularização
produziu uma alteração discreta no erro de validação. Como não foi
realizada uma busca pelo melhor valor de $\lambda_{L2}$, não é possível
generalizar esse comportamento para outras intensidades de
regularização.

\begin{figure}[H]
\centering
\includegraphics[width=0.90\textwidth]{figures/l2_random_curves.png}
\caption{Dinâmica de treinamento do modelo com regularização L2 no
particionamento aleatório.}
\label{fig:l2_random_curves}
\end{figure}

\begin{figure}[H]
\centering
\includegraphics[width=0.90\textwidth]{figures/l2_systematic_curves.png}
\caption{Dinâmica de treinamento do modelo com regularização L2 no
particionamento sistemático.}
\label{fig:l2_systematic_curves}
\end{figure}


\subsection{Baseline + Regularização L1}

A terceira ablação utilizou penalização pela norma L1 dos pesos, com
$\lambda_{L1}=10^{-4}$. A penalização foi aplicada aos pesos das
camadas, sem incluir os \textit{biases}.

\begin{table}[H]
\centering
\small
\caption{Resultados da ablação com regularização L1.}
\label{tab:l1_ablation}
\begin{tabular}{lrrrr}
\toprule
\textbf{Split} &
\textbf{MSE val.} &
\textbf{MAE val.} &
\textbf{Melhor época} &
\textbf{Épocas} \\
\midrule
Random &
0,805583 &
0,747984 &
623 &
643 \\

Systematic &
0,771047 &
0,734479 &
1000 &
1000 \\
\bottomrule
\end{tabular}
\end{table}

No particionamento aleatório, o L1 reduziu o MSE de validação de
$0,818838$ para $0,805583$, aproximadamente 1,62\%. A MAE passou de
$0,794267$ para $0,747984$.

No particionamento sistemático, o MSE foi $0,771047$, muito próximo ao
baseline de $0,770725$. A MAE foi $0,734479$.

Os resultados mostram que, com $\lambda_{L1}=10^{-4}$, a penalização
produziu alterações pequenas no erro de validação. A comparação com L2
deve ser interpretada dentro desse mesmo valor de regularização, pois
os hiperparâmetros não foram submetidos a uma nova busca para cada
método.

\begin{figure}[H]
\centering
\includegraphics[width=0.90\textwidth]{figures/l1_random_curves.png}
\caption{Dinâmica de treinamento do modelo com regularização L1 no
particionamento aleatório.}
\label{fig:l1_random_curves}
\end{figure}

\begin{figure}[H]
\centering
\includegraphics[width=0.90\textwidth]{figures/l1_systematic_curves.png}
\caption{Dinâmica de treinamento do modelo com regularização L1 no
particionamento sistemático.}
\label{fig:l1_systematic_curves}
\end{figure}


\subsection{Baseline + Dropout}

A quarta ablação consistiu na introdução de camadas Dropout com
probabilidade $p=0,2$ após as camadas ocultas.

\begin{table}[H]
\centering
\small
\caption{Resultados da ablação com Dropout.}
\label{tab:dropout_ablation}
\begin{tabular}{lrrrr}
\toprule
\textbf{Split} &
\textbf{MSE val.} &
\textbf{MAE val.} &
\textbf{Melhor época} &
\textbf{Épocas} \\
\midrule
Random &
1,020800 &
0,824734 &
39 &
59 \\

Systematic &
0,923539 &
0,769643 &
28 &
48 \\
\bottomrule
\end{tabular}
\end{table}

O Dropout produziu aumento do MSE de validação nos dois
particionamentos. No split aleatório, o MSE aumentou de $0,818838$
para $1,020800$, aproximadamente 24,66\%. No split sistemático, o
aumento foi de $0,770725$ para $0,923539$, aproximadamente 19,83\%.

Os melhores pontos de validação ocorreram nas épocas 39 e 28,
respectivamente. Portanto, embora o melhor ponto tenha ocorrido
precocemente, o erro associado a esse ponto foi maior que o observado
no baseline.

Uma possível interpretação é que, para a quantidade reduzida de
amostras de treinamento deste experimento, o Dropout com $p=0,2$ tenha
introduzido uma regularização excessiva para a configuração utilizada.
Essa hipótese é compatível com os resultados observados, mas é
específica do valor de $p$ testado e não permite concluir sobre outras
taxas de Dropout.

\begin{figure}[H]
\centering
\includegraphics[width=0.90\textwidth]{figures/dropout_random_curves.png}
\caption{Dinâmica de treinamento do modelo com Dropout no
particionamento aleatório.}
\label{fig:dropout_random_curves}
\end{figure}

\begin{figure}[H]
\centering
\includegraphics[width=0.90\textwidth]{figures/dropout_systematic_curves.png}
\caption{Dinâmica de treinamento do modelo com Dropout no
particionamento sistemático.}
\label{fig:dropout_systematic_curves}
\end{figure}


\subsection{Modelo Combinado}

Por fim, foi avaliada a combinação de momentum e regularização L2,
utilizando $\beta=0,9$ e $\lambda_{L2}=10^{-4}$. A combinação foi
introduzida para verificar se os efeitos observados individualmente
poderiam ser mantidos ou ampliados quando os dois mecanismos fossem
utilizados simultaneamente.

\begin{table}[H]
\centering
\small
\caption{Resultados do modelo combinado.}
\label{tab:combined_ablation}
\begin{tabular}{lrrrr}
\toprule
\textbf{Split} &
\textbf{MSE val.} &
\textbf{MAE val.} &
\textbf{Melhor época} &
\textbf{Épocas} \\
\midrule
Random &
0,723593 &
0,709413 &
128 &
148 \\

Systematic &
0,920469 &
0,769004 &
7 &
27 \\
\bottomrule
\end{tabular}
\end{table}

No particionamento aleatório, a combinação apresentou MSE de validação
$0,723593$, inferior ao baseline ($0,818838$), mas superior ao
resultado obtido com momentum isoladamente ($0,701855$). Dessa forma,
dentro desta configuração, a adição do L2 não produziu ganho adicional
em relação ao momentum.

No particionamento sistemático, o modelo combinado apresentou MSE de
validação $0,920469$, coincidindo com o valor observado para o modelo
com momentum isoladamente e permanecendo acima do baseline
($0,770725$).

Esses resultados mostram que os efeitos dos componentes não foram
aditivos de maneira direta. A combinação precisou ser tratada como uma
nova configuração experimental, e seu desempenho foi determinado pelo
erro de validação efetivamente observado.

\begin{figure}[H]
\centering
\includegraphics[width=0.90\textwidth]{figures/combined_random_curves.png}
\caption{Dinâmica de treinamento do modelo combinado no particionamento
aleatório.}
\label{fig:combined_random_curves}
\end{figure}

\begin{figure}[H]
\centering
\includegraphics[width=0.90\textwidth]{figures/combined_systematic_curves.png}
\caption{Dinâmica de treinamento do modelo combinado no particionamento
sistemático.}
\label{fig:combined_systematic_curves}
\end{figure}


\subsection{Comparação Consolidada das Ablações}

A Tabela~\ref{tab:ablation_comparison} consolida os resultados de
validação de todas as configurações avaliadas.

\begin{table}[H]
\centering
\small
\caption{Comparação consolidada das configurações avaliadas no conjunto
de validação.}
\label{tab:ablation_comparison}
\begin{tabular}{llrr}
\toprule
\textbf{Split} &
\textbf{Configuração} &
\textbf{MSE val.} &
\textbf{MAE val.} \\
\midrule
Random & Baseline & 0,818838 & 0,794267 \\
Random & Momentum & 0,701855 & 0,689057 \\
Random & L2 & 0,804006 & 0,747255 \\
Random & L1 & 0,805583 & 0,747984 \\
Random & Dropout & 1,020800 & 0,824734 \\
Random & Momentum + L2 & 0,723593 & 0,709413 \\
\midrule
Systematic & Baseline & 0,770725 & 0,734492 \\
Systematic & Momentum & 0,920469 & 0,769003 \\
Systematic & L2 & 0,770923 & 0,734615 \\
Systematic & L1 & 0,771047 & 0,734479 \\
Systematic & Dropout & 0,923539 & 0,769643 \\
Systematic & Momentum + L2 & 0,920469 & 0,769004 \\
\bottomrule
\end{tabular}
\end{table}

A seleção das configurações destinadas à avaliação final foi realizada
exclusivamente pelo critério definido de menor MSE de validação. Para
o particionamento aleatório, a configuração selecionada foi o modelo
baseline com momentum, com MSE de validação $0,701855$. Para o
particionamento sistemático, a configuração selecionada foi o modelo
baseline, com MSE de validação $0,770725$.

Essa seleção foi concluída antes da utilização do conjunto de teste,
que permaneceu separado das etapas de treinamento e escolha das
configurações.


\subsection{Avaliação Final no Conjunto de Teste}

Após a seleção baseada exclusivamente no conjunto de validação, cada
particionamento teve uma configuração definida para a avaliação final.
No particionamento aleatório, foi utilizado o modelo baseline com
momentum ($\beta=0,9$), enquanto no particionamento sistemático foi
utilizado o modelo baseline.

O conjunto de teste, composto por 240 observações em cada
particionamento, não foi utilizado durante o treinamento ou a seleção
dos modelos. Os resultados apresentados na
Tabela~\ref{tab:final_test} estão na escala original da variável
resposta.

\begin{table}[H]
\centering
\small
\caption{Desempenho final no conjunto de teste, no espaço original da
variável resposta.}
\label{tab:final_test}
\begin{tabular}{llrrrr}
\toprule
\textbf{Split} &
\textbf{Modelo} &
\textbf{MSE} &
\textbf{RMSE} &
\textbf{MAE} &
\textbf{$R^2$} \\
\midrule
Random &
Momentum &
0,447147 &
0,668690 &
0,556056 &
0,112666 \\

Systematic &
Baseline &
0,419493 &
0,647683 &
0,525382 &
0,190934 \\
\bottomrule
\end{tabular}
\end{table}

No particionamento aleatório, o modelo selecionado com momentum
apresentou MSE $0,447147$, RMSE $0,668690$, MAE $0,556056$ e
$R^2=0,112666$.

No particionamento sistemático, o modelo baseline apresentou MSE
$0,419493$, RMSE $0,647683$, MAE $0,525382$ e
$R^2=0,190934$.

As mesmas métricas calculadas no espaço padronizado foram:

\begin{table}[H]
\centering
\small
\caption{Desempenho final no conjunto de teste no espaço padronizado.}
\label{tab:final_test_scaled}
\begin{tabular}{llrrrr}
\toprule
\textbf{Split} &
\textbf{Modelo} &
\textbf{MSE} &
\textbf{RMSE} &
\textbf{MAE} &
\textbf{$R^2$} \\
\midrule
Random &
Momentum &
0,746377 &
0,863931 &
0,718410 &
0,112666 \\

Systematic &
Baseline &
0,725588 &
0,851815 &
0,690969 &
0,190934 \\
\bottomrule
\end{tabular}
\end{table}

O valor de $R^2$ não é alterado pela padronização da variável resposta,
enquanto MSE, RMSE e MAE dependem da escala utilizada. Por esse motivo,
os resultados finais são apresentados prioritariamente na escala
original, mantendo-se os valores padronizados para comparação com as
métricas utilizadas durante o treinamento.

A comparação entre validação e teste deve ser realizada separadamente
em cada particionamento, pois os conjuntos avaliados são distintos. No
split aleatório, o MSE passou de $0,701855$ na validação para
$0,746377$ no teste, um aumento de aproximadamente 6,34\%. No split
sistemático, o MSE passou de $0,770725$ na validação para $0,725588$ no
teste, uma redução de aproximadamente 5,86\%.

Essas diferenças não representam uma alteração do modelo entre as duas
etapas, mas sim uma diferença entre o desempenho observado em conjuntos
distintos. O conjunto de teste foi utilizado apenas após a seleção das
configurações, preservando sua função de avaliação final.


\subsection{Análise de Paridade e Resíduos}

Para complementar as métricas agregadas, foram analisadas as previsões
obtidas no conjunto de teste por meio de gráficos de paridade e de
resíduos.

Os gráficos de paridade comparam os valores reais e preditos, tendo
como referência a reta ideal

\[
\hat{y}=y.
\]

\begin{figure}[H]
\centering
\includegraphics[width=0.80\textwidth]{figures/parity_random.png}
\caption{Gráfico de paridade do modelo selecionado no particionamento
aleatório.}
\label{fig:parity_random}
\end{figure}

\begin{figure}[H]
\centering
\includegraphics[width=0.80\textwidth]{figures/parity_systematic.png}
\caption{Gráfico de paridade do modelo selecionado no particionamento
sistemático.}
\label{fig:parity_systematic}
\end{figure}

A análise dos resíduos foi realizada a partir de

\[
e_i=y_i-\hat{y}_i.
\]

As principais estatísticas descritivas dos resíduos no conjunto de
teste são apresentadas na Tabela~\ref{tab:residual_summary}.

\begin{table}[H]
\centering
\scriptsize
\caption{Resumo estatístico dos resíduos no conjunto de teste.}
\label{tab:residual_summary}
\begin{tabular}{lrrrrrrrrrr}
\toprule
\textbf{Split} &
\textbf{Média} &
\textbf{Mediana} &
\textbf{DP} &
\textbf{$Q_1$} &
\textbf{$Q_3$} &
\textbf{IQR} &
\textbf{Mín.} &
\textbf{Máx.} &
\textbf{Skew.} &
\textbf{Curt.} \\
\midrule
Random
& -0,091612
& -0,129840
& 0,662385
& -0,578163
& 0,445578
& 1,023741
& -1,811030
& 2,277711
& 0,147780
& -0,153236 \\

Systematic
& 0,069106
& -0,004271
& 0,643986
& -0,382244
& 0,522215
& 0,904458
& -1,780099
& 1,986360
& 0,184936
& -0,217566 \\
\bottomrule
\end{tabular}
\end{table}

No particionamento aleatório, a média dos resíduos foi $-0,091612$ e a
mediana foi $-0,129840$. Como $e=y-\hat{y}$, resíduos negativos
correspondem a previsões superiores aos valores observados. A
combinação de média e mediana negativas indica uma pequena tendência
nessa direção.

No particionamento sistemático, a média foi $0,069106$, enquanto a
mediana foi $-0,004271$. A mediana permanece muito próxima de zero,
enquanto a média positiva mostra a influência de parte dos resíduos
positivos sobre a distribuição.

A dispersão central dos resíduos foi semelhante nos dois
particionamentos. O intervalo interquartil foi $1,023741$ no random e
$0,904458$ no systematic. Os valores extremos variaram de
$-1,811030$ a $2,277711$ no random e de $-1,780099$ a $1,986360$ no
systematic.

As medidas de assimetria foram pequenas e positivas: $0,147780$ no
random e $0,184936$ no systematic. Dessa forma, há uma assimetria
positiva leve, sem evidência de assimetria acentuada. A curtose em
excesso foi $-0,153236$ no random e $-0,217566$ no systematic. Esses
valores descrevem a forma da distribuição dos resíduos, mas não
constituem, isoladamente, um teste de normalidade.

A relação entre erro médio e erro quadrático pode ser expressa por

\[
\operatorname{MSE}
=
\operatorname{Var}(e)
+
[\operatorname{E}(e)]^2.
\]

Para o particionamento aleatório, o termo associado à média dos
resíduos corresponde a aproximadamente 1,88\% do MSE. No
particionamento sistemático, essa parcela corresponde a aproximadamente
1,14\%. Portanto, a maior parte do erro quadrático observado está
associada à dispersão dos resíduos, e não ao seu deslocamento médio.

Os resíduos em função de $x$ são apresentados nas
Figuras~\ref{fig:residual_x_random} e
\ref{fig:residual_x_systematic}. Essa representação é particularmente
relevante neste problema, pois $x$ é a única variável de entrada e a
relação observada entre $x$ e $y$ é não linear.

\begin{figure}[H]
\centering
\includegraphics[width=0.90\textwidth]{figures/residuals_x_random.png}
\caption{Resíduos em função de $x$ no particionamento aleatório.}
\label{fig:residual_x_random}
\end{figure}

\begin{figure}[H]
\centering
\includegraphics[width=0.90\textwidth]{figures/residuals_x_systematic.png}
\caption{Resíduos em função de $x$ no particionamento sistemático.}
\label{fig:residual_x_systematic}
\end{figure}

Finalmente, os histogramas dos resíduos são apresentados nas
Figuras~\ref{fig:residual_hist_random} e
\ref{fig:residual_hist_systematic}. Essas representações complementam
as estatísticas descritivas ao permitir a inspeção visual da
concentração dos erros e de suas caudas.

\begin{figure}[H]
\centering
\includegraphics[width=0.90\textwidth]{figures/residual_hist_random.png}
\caption{Distribuição dos resíduos no particionamento aleatório.}
\label{fig:residual_hist_random}
\end{figure}

\begin{figure}[H]
\centering
\includegraphics[width=0.90\textwidth]{figures/residual_hist_systematic.png}
\caption{Distribuição dos resíduos no particionamento sistemático.}
\label{fig:residual_hist_systematic}
\end{figure}


\section{Discussão e Lições Aprendidas}

Os experimentos mostraram que, no regime utilizado, o desempenho da
MLP depende não apenas da arquitetura escolhida, mas também da
composição das partições e do mecanismo de treinamento empregado.

O \texttt{dataset\_p1} possui 300 observações, mas somente 30 foram
utilizadas para treinamento em cada estratégia. Assim, o modelo foi
ajustado com uma quantidade reduzida de exemplos, enquanto 240
observações permaneceram separadas para a avaliação final.

\subsection{Efeito do Regime de Baixa Disponibilidade de Dados}

A utilização de apenas 10\% das observações para treinamento torna a
estimativa da relação entre $x$ e $y$ dependente da composição das
amostras utilizadas para ajuste e validação. A diferença entre os dois
particionamentos evidencia esse efeito.

No particionamento aleatório, o uso de momentum reduziu o MSE de
validação em aproximadamente 14,29\%, enquanto L1 e L2 produziram
reduções menores, de aproximadamente 1,62\% e 1,81\%. O Dropout com
$p=0,2$, por sua vez, aumentou o MSE em aproximadamente 24,66\%.

No particionamento sistemático, o baseline apresentou MSE de validação
$0,770725$. L1 e L2 permaneceram muito próximos desse valor, enquanto
momentum e Dropout elevaram o MSE para $0,920469$ e $0,923539$,
respectivamente.

A diferença entre esses comportamentos mostra que os resultados de uma
ablação não podem ser dissociados da forma como os dados foram
particionados. Em particular, utilizar apenas uma divisão poderia
atribuir a uma técnica um efeito que não se reproduz sob uma
partição diferente.


\subsection{Comportamento do SGD e Efeito do Momentum}

O baseline foi treinado com SGD sem momentum. No particionamento
aleatório, a introdução de $\beta=0,9$ reduziu o MSE de validação de
$0,818838$ para $0,701855$ e deslocou a melhor época de 564 para 152.

No particionamento sistemático, o mesmo mecanismo aumentou o MSE de
$0,770725$ para $0,920469$, enquanto o melhor ponto ocorreu na época 7.

Assim, o experimento evidencia que alterar a dinâmica do otimizador
pode produzir efeitos dependentes da distribuição das observações entre
treinamento e validação. A diferença também reforça que a época de
convergência não deve ser interpretada isoladamente: no systematic, o
modelo com momentum atingiu rapidamente seu melhor ponto, mas esse
ponto apresentou erro de validação superior ao baseline.


\subsection{Comparação entre L1 e L2}

As regularizações L1 e L2 foram avaliadas com $\lambda=10^{-4}$. No
particionamento aleatório, os MSEs foram $0,805583$ para L1 e
$0,804006$ para L2, ambos abaixo do baseline de $0,818838$.

No particionamento sistemático, L1 apresentou MSE $0,771047$ e L2
apresentou $0,770923$, enquanto o baseline apresentou $0,770725$.

Os resultados indicam que, para a intensidade de regularização
utilizada, os efeitos de L1 e L2 sobre o MSE de validação foram
pequenos. Não foi realizada uma busca pelos valores de $\lambda$ para
essas regularizações. Portanto, a comparação diz respeito às
configurações testadas e não à capacidade máxima de cada método sob uma
otimização específica de seu hiperparâmetro.


\subsection{Efeito do Dropout}

O Dropout com $p=0,2$ aumentou o erro de validação nos dois
particionamentos. O MSE passou de $0,818838$ para $1,020800$ no random
e de $0,770725$ para $0,923539$ no systematic.

Os melhores pontos ocorreram nas épocas 39 e 28, respectivamente.
Portanto, nesse experimento, o treinamento atingiu o melhor ponto
rapidamente, mas o ponto obtido apresentou erro de validação maior que
o baseline.

Uma interpretação possível é que a aleatoriedade introduzida pelo
Dropout tenha representado uma regularização excessiva diante das
apenas 30 amostras disponíveis para treinamento. Essa interpretação é
compatível com os resultados, mas permanece limitada ao valor
$p=0,2$ utilizado.


\subsection{Interação entre Momentum e L2}

A combinação entre momentum e L2 apresentou MSE de validação $0,723593$
no random e $0,920469$ no systematic.

No random, a combinação permaneceu melhor que o baseline, mas não
superou o momentum isolado. No systematic, o resultado coincidiu com
o experimento de momentum e permaneceu acima do baseline.

Portanto, nesta configuração, os efeitos individuais não foram
simplesmente aditivos. A regularização L2 não acrescentou uma redução
adicional do erro ao momentum e, em nenhum dos dois particionamentos,
a configuração combinada superou simultaneamente todas as configurações
individuais.


\subsection{Generalização para o Conjunto de Teste}

A seleção dos modelos foi realizada exclusivamente com base no conjunto
de validação. O conjunto de teste permaneceu reservado para a avaliação
final, evitando que seu desempenho influenciasse a escolha da
configuração.

No particionamento aleatório, o modelo selecionado com momentum
apresentou MSE de teste $0,447147$, RMSE $0,668690$, MAE $0,556056$ e
$R^2=0,112666$.

No particionamento sistemático, o modelo baseline apresentou MSE de
teste $0,419493$, RMSE $0,647683$, MAE $0,525382$ e
$R^2=0,190934$.

No random, o MSE de teste foi aproximadamente 6,34\% maior que o MSE de
validação. No systematic, o MSE de teste foi aproximadamente 5,86\%
menor que o MSE de validação.

Essas diferenças não permitem inferir, isoladamente, que um conjunto
seja mais difícil que o outro ou que o modelo tenha mudado de
comportamento entre as etapas. Elas são resultado da avaliação em
amostras distintas e devem ser interpretadas dentro de cada
particionamento.


\subsection{Análise dos Resíduos}

Os resíduos apresentaram dispersão semelhante nos dois modelos
selecionados. O desvio-padrão foi $0,662385$ no random e $0,643986$ no
systematic.

No random, a média e a mediana foram negativas, respectivamente
$-0,091612$ e $-0,129840$, indicando uma pequena tendência média de
superestimação. No systematic, a média foi $0,069106$ e a mediana foi
$-0,004271$, o que indica que a posição central da distribuição ficou
muito próxima de zero, apesar da influência de resíduos positivos sobre
a média.

A skewness positiva e de pequena magnitude nos dois particionamentos é
compatível com uma assimetria leve. Os valores extremos também mostram
que existem observações com erros significativamente maiores que a
magnitude típica dos resíduos. Como o MSE atribui peso quadrático a
esses desvios, algumas poucas observações podem contribuir de maneira
relevante para o erro agregado.

A decomposição do MSE mostrou ainda que apenas uma pequena parcela do
erro quadrático está associada ao deslocamento médio dos resíduos.
Portanto, a principal contribuição para o erro observado está
associada à variabilidade dos desvios individuais.

A análise conjunta das métricas e dos gráficos de resíduos é, assim,
mais informativa que qualquer uma dessas medidas isoladamente.


\subsection{Lições Aprendidas}

O experimento evidenciou a importância de separar as etapas de
otimização, seleção e avaliação final. O conjunto de teste não foi
utilizado para selecionar a arquitetura ou os componentes do modelo,
sendo reservado para a estimativa final de generalização.

A busca inicial mostrou que a configuração do baseline também depende
da estratégia de particionamento, embora ambas tenham selecionado a
mesma arquitetura e ativação. Ao longo das ablações, os efeitos das
técnicas foram diferentes entre as partições: momentum apresentou
redução importante do MSE no random, mas aumento no systematic; L1 e L2
produziram alterações pequenas; e Dropout com $p=0,2$ aumentou o erro
nos dois casos.

A combinação entre momentum e L2 também não apresentou ganho adicional
em relação ao momentum isolado no random e não superou o baseline no
systematic. Isso mostra que a combinação de mecanismos de otimização e
regularização precisa ser verificada experimentalmente, não sendo
suficiente supor que seus efeitos serão cumulativos.

Por fim, a análise do conjunto de teste mostrou que a avaliação por
MSE, RMSE, MAE e $R^2$ deve ser complementada pela inspeção das curvas
de treinamento, dos gráficos de paridade e dos resíduos. Essa análise
conjunta permite distinguir o comportamento de otimização da capacidade
de generalização e explicitar as limitações do experimento, sobretudo
diante da quantidade reduzida de observações utilizadas para
treinamento.

\end{document}